\documentclass[10pt,letterpaper]{article}
\usepackage{cornell-notes}

\title{Clause 26.07.28.02: Class Template Chunk View For Input Ranges}
\author{}
\date{}
\setDocTitle{Clause 26.07.28.02: Class Template Chunk View For Input Ranges}
\setDocSubtitle{C++ 2024}
\setDocOwner{C++ 2024 Cornell Notes}
\setDocDate{\today}
\setCornellCollection{C++ 2024 Cornell Notes}
\setCornellUnitType{Clause}
\setCornellUnitNumber{26.07.28.02}
\setCornellUnitTitle{Class Template Chunk View For Input Ranges}
\hypersetup{pdftitle={Clause 26.07.28.02: Class Template Chunk View For Input Ranges},pdfsubject={Cornell notes},pdfauthor={}}

\begin{document}
\maketitle
\begin{CornellOverviewBox}{Overview}
\textbf{Classification:} Standard library - Normative\\[2pt]
\textbf{Learning objective:} Understand the rules, terminology, and semantic consequences associated with Class template chunk\_view for input ranges.
\end{CornellOverviewBox}\begin{CornellSummaryBox}{Summary}
{\sffamily\bfseries\color{Accent} Summary}\par\vspace{3pt}Section 26.7.28.2 focuses on Class template chunk\_view for input ranges. Apply the specified interface together with its mandates, effects, result conditions, exception behavior, and complexity constraints. When applying it, preserve the distinction between mandatory requirements, recommendations, permissions, and behavior deliberately left undefined, unspecified, or implementation-defined.
\end{CornellSummaryBox}

\end{document}
